\section{Information Extraction}\label{sec:information_extraction_evaluation}

The information extraction module extracts four key attributes from the product pages: the product's title, price, image URL, and availability status.
Since all four of these attributes are independent, their extraction performance can be evaluated separately.
As mentioned in section~\ref{subsec:information_extraction_training_phase}, this is fundamentally a token classification task, where each token in the HTML page is classified as either belonging to one of the four attributes or not.
Therefore, we can track the F1-score for each of the four attributes separately, as well as the overall F1-score for the information extraction module.

Table~\ref{tab:information-extraction-classification-report} shows the precision, recall, and F1-score for each label, along with the macro average.
The overall F1-score for the information extraction module is 0.9700, which indicates that the module is able to accurately extract the required attributes from the product pages.

\begin{table}[H]
    \centering
    \begin{tabular}{lrrr}
        \toprule
        \textbf{Label} & \textbf{Precision} & \textbf{Recall} & \textbf{F1} \\
        \midrule
        O & 0.9953 & 0.9931 & 0.9942 \\
        B-TITLE & 0.9408 & 0.9420 & 0.9414 \\
        B-PRICE & 0.9719 & 0.9745 & 0.9732 \\
        B-IMAGE & 0.9310 & 0.9776 & 0.9537 \\
        B-OUT\_OF\_STOCK & 0.9798 & 0.9949 & 0.9873 \\
        \midrule
        \textbf{MACRO} & \textbf{0.9638} & \textbf{0.9764} & \textbf{0.9700} \\
        \bottomrule
    \end{tabular}
    \caption{Information extraction label-wise classification performance}
    \label{tab:information-extraction-classification-report}
\end{table}

These figures may deviate slightly from the results reported in section~\ref{subsec:information_extraction_model_evaluation} due to the use of a larger training dataset and a different validation dataset for the evaluation of the information extraction module.

As with the page classifier module, the information extraction module was also evaluated for its inference speed.
Three of the last runs of the information extraction module were timed, and the time taken for each phase was recorded.
Table~\ref{tab:information-extraction-runtime} summarizes the runtime across three runs.

\begin{table}[H]
    \centering
    \begin{tabular}{lrrrrr}
        \toprule
        \textbf{Run} & \textbf{Pages} & \textbf{Generate Dataset (s)} & \textbf{Inference (s)} & \textbf{Update Database (s)} & \textbf{Total (s)} \\
        \midrule
        1 & 26,010 & 570.30 & 1,501.08 & 104.00 & 2,175.38 \\
        2 & 26,093 & 610.52 & 1,566.50 & 98.84 & 2,275.86 \\
        3 & 27,148 & 660.53 & 1,571.18 & 98.79 & 2,330.51 \\
        \midrule
        \textbf{Total} & \textbf{79,251} & \textbf{1,841.35} & \textbf{4,638.76} & \textbf{301.63} & \textbf{6,781.74} \\
        \bottomrule
    \end{tabular}
    \caption{Information extraction runtime by phase across runs}
    \label{tab:information-extraction-runtime}
\end{table}

A total of \texttt{79,251} product pages were processed across the three runs, with a total runtime of \texttt{6,781.74} seconds.
Therefore, for a batch of \texttt{20,000} product pages, the average runtime is \texttt{1,710.00} seconds, which is approximately \texttt{28.5} minutes.